\section{Simulation by an ordinary bimodule}
\label{sec:ordinary-simulation}

We now construct a larger finite-dimensional algebra and an ordinary
bimodule whose second derived tensor iterate recovers \(P\otL_B-\),
up to a fixed shift. All choices in this section depend only on the
fixed complex \(P\) of \Cref{thm:realization}.

Choose integers \(a_*\le 0\le b_*\) such that \(P^n=0\) outside
\([a_*,b_*]\), enlarging the interval if necessary, and put
\(l=b_*-a_*\).
Let \(C\) be the algebra of the quiver
\[
0\xrightarrow{c_1}1\xrightarrow{c_2}\cdots
\xrightarrow{c_l}l
\]
modulo all paths of length two.
Write \(\varepsilon_i\) for its vertex idempotents, and let \(W\)
be the simple \emph{right} \(C\)-module at vertex \(l\).
There is a projective resolution of right \(C\)-modules
\begin{equation}
\label{eq:right-chain-resolution}
0\longrightarrow \varepsilon_0C
\xrightarrow{c_1\cdot}\varepsilon_1C
\xrightarrow{c_2\cdot}\cdots
\xrightarrow{c_l\cdot}\varepsilon_lC
\longrightarrow W\longrightarrow0.
\end{equation}
Indeed, for \(i>0\), the right projective \(\varepsilon_iC\)
has basis \(\varepsilon_i,c_i\), and its radical is the right
simple at \(i-1\).  Left multiplication by \(c_i\) maps
\(\varepsilon_{i-1}\) to \(c_i\) and, when \(i>1\), maps
\(c_{i-1}\) to zero.  This proves exactness, including injectivity
at the first term.  For \(l=0\), the resolution consists of
the isomorphism \(\varepsilon_0C\to W\).

Define a \((B\otimes_k C,B)\)-bimodule \(O\) by
\[
\varepsilon_iO=P^{a_*+i},
\qquad
c_i\big|_{\varepsilon_{i-1}O}=d_P^{a_*+i-1}.
\]
Here the two \(B\)-actions are those on the terms of \(P\).
The relation \(d_P^2=0\) makes this a representation of \(C\),
and the bimodule linearity of \(d_P\) makes the \(C\)-action
commute with both \(B\)-actions.  Thus \(O\) is an ordinary
finite-dimensional bimodule.  Set
\[
D=B\times(B\otimes_k C),\qquad
T=B\otimes_k W,\qquad X=O\oplus T.
\]
Regard \(O\) as having its left action from the second factor
of \(D\) and its right action from the first.  Regard \(T\)
as having its left action from the first factor and its
right action from the second, with the regular \(B\)-actions
on its \(B\)-factor.  These prescriptions define an ordinary
finite-dimensional \(D\)-bimodule \(X\).

The two summands of \(X\) move complexes between the factors of \(D\),
as shown in \Cref{fig:simulation}. Thus two tensor steps return to the
first factor; the next proposition identifies their composite.

\begin{figure}[htbp]
\centering
\begin{tikzcd}[column sep=huge]
 \D^b(B\text{-}\modu)
 \arrow[r,bend left=22,"O\otL_B-"] &
 \D^b((B\otimes_k C)\text{-}\modu)
 \arrow[l,bend left=22,"T\otL_{B\otimes_k C}-"]
\end{tikzcd}
\caption{The two factors of \(D\) and the tensor functors supplied by
the off-diagonal summands of \(X\). Starting on the left, the composite
is \(P[b_*]\otL_B-\) by \Cref{prop:ordinary-simulation}.}
\label{fig:simulation}
\end{figure}

\begin{proposition}
\label{prop:ordinary-simulation}
The algebra \(D\) has finite global dimension on both sides.
Let \(F=X\otL_D-\).  For every complex \(Z\) of left
\(B\)-modules, regarded as a complex on the first factor of
\(D\), there is an isomorphism
\begin{equation}
\label{eq:compression}
F^2(Z)\simeq P[b_*]\otL_B Z,
\qquad
T\otL_{B\otimes_k C}O\simeq P[b_*].
\end{equation}
The second isomorphism is an isomorphism of complexes of
\(B\)-bimodules, and the first is supported on the first
factor of \(D\).
\end{proposition}

\begin{proof}
The vertices of \(B\otimes_k C\) are pairs \((i,j)\), where
\(0\le i\le2\) and \(0\le j\le l\).  Its diagonal vertex
corners are \(k\).  A nonzero component between distinct
vertices moves in the product of the two directed orders.
Choosing a linear extension of that product order permits
an application of \Cref{lem:directed-bound}, on both sides.
In particular,
\[
\gldim(B\otimes_k C)\le3(l+1)-1,
\qquad
\gldim\bigl((B\otimes_k C)^{\mathrm{op}}\bigr)
\le3(l+1)-1.
\]
\Cref{lem:directed-bound} also gives the bound \(2\) for \(B\) and
\(B^{\mathrm{op}}\).  Modules over a finite product split
according to its factors, so
\[
\gldim D,\ \gldim D^{\mathrm{op}}\le3l+2.
\]

Tensor \eqref{eq:right-chain-resolution} over \(k\) with
the regular \(B\)-bimodule, placing its last projective term
in degree zero.  This gives a bounded resolution \(Q_T\to T\)
by right \(B\otimes_k C\)-projectives, preserving the left
\(B\)-action.  Also, \(O\) is projective as a right
\(B\)-module, being the finite direct sum of the terms of \(P\).
Thus \(Q_X=O\oplus Q_T\), with the prescribed block actions,
is a bounded right-\(D\)-projective resolution of \(X\) retaining
its left \(D\)-action.  Ordinary tensor with \(Q_X\) computes
\(F\) on every complex.  On the first factor, the block supports
therefore give
\[
F^2(Z)\simeq
Q_T\otimes_{B\otimes_k C}(O\otimes_B Z)
\cong
\bigl(Q_T\otimes_{B\otimes_k C}O\bigr)\otimes_B Z.
\]
In degree \(n\), for \(-l\le n\le0\), tensoring \(Q_T\)
with \(O\) gives
\[
\bigl(B\otimes_k\varepsilon_{l+n}C\bigr)
\otimes_{B\otimes_k C}O
\cong\varepsilon_{l+n}O
=P^{a_*+l+n}=P^{b_*+n}.
\]
The induced differential is \(d_P\).
In comparison, the standard differential of \(P[b_*]\) is
\((-1)^{b_*}d_P\).  Multiplication by
\((-1)^{b_*n}\) in degree \(n\) defines a chain isomorphism
to \(P[b_*]\); it respects both \(B\)-actions.
Since this complex is termwise right-\(B\)-projective, its
ordinary tensor with any \(Z\) computes the derived tensor.
This proves the second isomorphism in
\eqref{eq:compression}, and hence the first.
\end{proof}
